%=============================================================================!
% 5.5  THE INERTIA TENSOR                                                     !
%=============================================================================!
\section{The inertia tensor}
\label{sec:ro-tensor}

\Cref{sec:ro-rigid} found the angular momentum along the axis of turning,
$L_z = I\omega$. A body can also have angular momentum across that axis:
two weights fixed to the ends of a rod, the rod clamped at a slant to an
axle, keep pulling sideways on the axle as they go round. This section
finds the whole vector $\vb{L}$ for a rigid body turning in any direction,
and with it the matrix that takes the place of the mass.

\subsection*{Angular momentum of a turning body}

Let a rigid body turn with angular velocity $\boldsymbol{\omega}$ about an
axis through its centre of mass, and measure each position $\vb{r}_i$ from
the centre of mass. Each part moves with $\vb{v}_i = \boldsymbol{\omega}
\times\vb{r}_i$, so
\begin{equation*}
   \vb{L} = \sum_i m_i\,\vb{r}_i\times(\boldsymbol{\omega}\times\vb{r}_i) .
\end{equation*}
Two rules for products of three vectors turn this into something
computable.

\begin{toolbox}[Triple products]
\emph{The vector triple product.} For any three vectors,
\begin{equation*}
   \vb{a}\times(\vb{b}\times\vb{c})
   = \vb{b}\,(\vb{a}\cdot\vb{c}) - \vb{c}\,(\vb{a}\cdot\vb{b}) .
\end{equation*}
To check the $x$ component, use $(\vb{b}\times\vb{c})_z = b_xc_y -
b_yc_x$ and $(\vb{b}\times\vb{c})_y = b_zc_x - b_xc_z$:
\begin{align*}
   \bigl[\vb{a}\times(\vb{b}\times\vb{c})\bigr]_x
   &= a_y(b_xc_y - b_yc_x) - a_z(b_zc_x - b_xc_z)\\
   &= b_x(a_yc_y + a_zc_z) - c_x(a_yb_y + a_zb_z)
   && \text{multiplying out}\\
   &= b_x(\vb{a}\cdot\vb{c}) - c_x(\vb{a}\cdot\vb{b})
   && \text{adding and subtracting $a_xb_xc_x$.}
\end{align*}
The $y$ and $z$ components follow by renaming $x \to y \to z \to x$.

\emph{The scalar triple product.} Writing out
\begin{equation*}
   \vb{a}\cdot(\vb{b}\times\vb{c})
   = a_x(b_yc_z - b_zc_y) + a_y(b_zc_x - b_xc_z) + a_z(b_xc_y - b_yc_x)
\end{equation*}
gives six products, each of one component of every vector. Writing out
\begin{equation*}
   \vb{b}\cdot(\vb{c}\times\vb{a})
   = b_x(c_ya_z - c_za_y) + b_y(c_za_x - c_xa_z) + b_z(c_xa_y - c_ya_x)
\end{equation*}
in the same way gives the same six with
the same signs, in a different order. So the three vectors may be moved
round in a cycle, $\vb{a}\cdot(\vb{b}\times\vb{c}) = \vb{b}\cdot(\vb{c}
\times\vb{a}) = \vb{c}\cdot(\vb{a}\times\vb{b})$.
\end{toolbox}

By the vector triple product with $\vb{a} = \vb{c} = \vb{r}_i$ and $\vb{b}
= \boldsymbol{\omega}$, $\vb{r}_i\times(\boldsymbol{\omega}\times\vb{r}_i)
= \boldsymbol{\omega}\,\lvert\vb{r}_i\rvert^2 - \vb{r}_i\,(\vb{r}_i\cdot
\boldsymbol{\omega})$. Its component $a$, for $a = x$, $y$ or $z$, is
$\omega_a\lvert\vb{r}_i\rvert^2 - (\vb{r}_i)_a\sum_b(\vb{r}_i)_b\omega_b$,
with $(\vb{r}_i)_a$ the component $a$ of $\vb{r}_i$. Writing $\omega_a = \sum_b\delta_{ab}
\omega_b$, where the Kronecker delta $\delta_{ab}$ is 1 when $a = b$ and 0
otherwise, so that only the term with $b = a$ survives, both terms become
sums over $b$; multiplying by $m_i$, adding over the atoms and taking the
sum over $b$ outside,
\begin{equation}
   L_a = \sum_b I_{ab}\,\omega_b , \qquad
   I_{ab} = \sum_i m_i\,\bigl(\lvert\vb{r}_i\rvert^2\delta_{ab}
   - (\vb{r}_i)_a(\vb{r}_i)_b\bigr) ,
   \label{eq:ro-tensor}
\end{equation}
or $\vb{L} = \mathsf{I}\boldsymbol{\omega}$ with the matrix product of
\cref{sec:os-modes}. The $3\times3$ matrix $\mathsf{I}$, set in sans-serif
as this book sets tensors, apart from the identity $\vb{I}$, is the
inertia tensor of the body about its centre of mass. It is symmetric,
$I_{ab} = I_{ba}$, since swapping $a$ and $b$ changes neither
$\delta_{ab}$ nor $(\vb{r}_i)_a(\vb{r}_i)_b$. Its diagonal entries are the moments of inertia about the three
axes, for instance $I_{zz} = \sum_i m_i(x_i^2 + y_i^2 + z_i^2 - z_i^2) =
\sum_i m_i(x_i^2 + y_i^2)$, as in \cref{eq:ro-moment}; its other entries,
such as $I_{xy} = -\sum_i m_ix_iy_i$, are the source of the sideways
angular momentum. For $\boldsymbol{\omega} = (0, 0, \omega)$, $\vb{L} =
\omega\,(I_{xz}, I_{yz}, I_{zz})$, along the axis only when $I_{xz}$ and
$I_{yz}$ vanish.

The kinetic energy follows with the scalar triple product. Since $\vb{v}_i
= \boldsymbol{\omega}\times\vb{r}_i$,
\begin{equation*}
   \lvert\vb{v}_i\rvert^2 = \vb{v}_i\cdot(\boldsymbol{\omega}\times\vb{r}_i)
   = \boldsymbol{\omega}\cdot(\vb{r}_i\times\vb{v}_i) ,
\end{equation*}
moving the vectors round in a cycle, so that multiplying by $\frac12 m_i$
and adding, with $\sum_i m_i\,\vb{r}_i\times\vb{v}_i = \vb{L}$,
\begin{equation}
   K = \tfrac12\,\boldsymbol{\omega}\cdot\vb{L}
   = \tfrac12\,\boldsymbol{\omega}\cdot\mathsf{I}\boldsymbol{\omega} .
   \label{eq:ro-kinetic}
\end{equation}

\subsection*{Principal axes}

Because $\mathsf{I}$ is symmetric, it has three real eigenvalues $I_1 \le
I_2 \le I_3$ with eigenvectors at right angles to one another
(\cref{sec:os-modes}). The eigenvectors are the principal axes of the
body and the eigenvalues its principal moments of inertia. Turning about
a principal axis, $\mathsf{I}\boldsymbol{\omega} = I_k\boldsymbol{\omega}$,
and the angular momentum lies along the axis; turning about any other axis,
it does not. Write $\boldsymbol{\omega} = \omega_1\vb{e}_1 + \omega_2\vb{e}_2 +
\omega_3\vb{e}_3$, with $\vb{e}_k$ the principal axes as vectors of length
1 and $\omega_k$ the parts of $\boldsymbol{\omega}$ along them. The product
is linear, so $\mathsf{I}\boldsymbol{\omega} = I_1\omega_1\vb{e}_1 +
I_2\omega_2\vb{e}_2 + I_3\omega_3\vb{e}_3$, and since $\vb{e}_j\cdot
\vb{e}_k$ is 1 for $j = k$ and 0 otherwise, \cref{eq:ro-kinetic} becomes
$K = \frac12(I_1\omega_1^2 + I_2\omega_2^2 + I_3\omega_3^2)$, as for the
quadratic term in \cref{sec:os-modes}. With the principal axes as the
coordinate axes, $\vb{e}_1 = (1, 0, 0)$ and so on, and column $k$ of the
tensor, $\mathsf{I}\vb{e}_k = I_k\vb{e}_k$, holds $I_k$ on the diagonal
and zeros elsewhere: the tensor is diagonal.

A wheel turning steadily about an axle that is not a principal axis has
an angular momentum that goes round with the wheel: worked out in axes
fixed in the wheel, with the axle as their $z$ axis, $\vb{L} =
\omega\,(I_{xz}, I_{yz}, I_{zz})$ does not change, so its sideways part
turns with the wheel. Its direction changes, so by \cref{eq:ro-system}
the bearings must supply an external torque all the time, and the axle
shakes. Car wheels are balanced by clipping small weights to
their rims until the axle passes through the centre of mass and is a
principal axis.

\subsection*{A water molecule}

The measured equilibrium geometry of water\cite{CCCBDB} has O--H bonds of
\SI{0.958}{\angstrom} at an angle of \ang{104.4776}. Place the oxygen atom
at the origin, one hydrogen atom on the $x$ axis, at $(0.958, 0, 0)$, and
the other in the $xy$ plane, at
\begin{equation*}
   (0.958\cos\ang{104.4776},\ 0.958\sin\ang{104.4776},\ 0)
   = (-0.2395, 0.9276, 0)\ \si{\angstrom}
\end{equation*}
(\cref{fig:ro-water}), with the masses 15.999 and \SI{1.008}{\amu}. The
centre of mass is at $1.008\times(0.7185, 0.9276, 0)/18.015 = (0.0402,
0.0519, 0)$\,\si{\angstrom}, and about it
\cref{eq:ro-tensor} gives, in \si{\amu\angstrom\squared},
\begin{equation*}
   \mathsf{I} = \begin{pmatrix}
      0.8188 & 0.2615 & 0\\
      0.2615 & 0.9538 & 0\\
      0 & 0 & 1.7726
   \end{pmatrix} .
\end{equation*}
The atoms lie in the plane $z = 0$, so $I_{xz} = I_{yz} = 0$ and the $z$
axis, at right angles to the molecule, is a principal axis, with $I_3 =
1.7726$. The other two come from the $2\times2$ block in the top left, by
the characteristic equation of \cref{sec:os-modes}, multiplied out with
$0.8188 + 0.9538 = 1.7726$ and $0.8188\times0.9538 - 0.2615^2 = 0.7126$:
\begin{equation*}
   (0.8188 - \lambda)(0.9538 - \lambda) - 0.2615^2
   = \lambda^2 - 1.7726\,\lambda + 0.7126 = 0 ,
\end{equation*}
so $\lambda = \frac12\bigl(1.7726 \pm\sqrt{1.7726^2 - 4\times0.7126}\bigr)
= \frac12(1.7726 \pm 0.5401)$, giving $I_1 = 0.616$ and $I_2 =
\SI{1.156}{\amu\angstrom\squared}$. The axis of $I_2$ is the line that
halves the angle between the bonds, an axis of symmetry of the molecule:
putting $\lambda = 1.1564$ into the first equation of the pair in
\cref{sec:os-modes}, $(0.8188 - 1.1564)\,u_1 + 0.2615\,u_2 = 0$, gives
$u_2/u_1 = 0.3376/0.2615 = 1.291$, a line at $\arctan1.291 = \ang{52.24}$
to the $x$ axis, half of \ang{104.4776}. The axis of $I_1$ is at right
angles to it in the plane. The two roots, $\frac12(1.7726 \pm 0.5401)$,
add to $1.7726 = 0.8188 + 0.9538 = I_{xx} + I_{yy}$, and for a flat body
$I_{xx} + I_{yy} = I_{zz}$ (\cref{ex:ro-flat}), so the largest moment, neither of the other two being negative, each a
sum of masses times squared distances, is
the sum of the other two; computed without rounding, $I_1 = 0.6162$ and
$I_2 = 1.1564$.

\begin{figure}[htb]
\centering
\includegraphics{ch05_rotation/water}
\caption{A water molecule at its measured geometry, one bond along $x$,
with its centre of mass (cross) and the two principal axes in its plane,
labelled with their principal moments in \si{\amu\angstrom\squared}. The
third principal axis points out of the page, with $I_3 =
\SI{1.773}{\amu\angstrom\squared}$.}
\label{fig:ro-water}
\end{figure}

A linear molecule, with every atom on one line, has a zero principal
moment about that line, since every $\rho_i$ is zero, and two equal ones
about the axes at right angles to it.

\begin{keyidea}
For a rigid body turning about its centre of mass, $\vb{L} = \mathsf{I}
\boldsymbol{\omega}$ and $K = \frac12\boldsymbol{\omega}\cdot\mathsf{I}
\boldsymbol{\omega}$, with the symmetric inertia tensor $I_{ab} = \sum_i
m_i(\lvert\vb{r}_i\rvert^2\delta_{ab} - (\vb{r}_i)_a(\vb{r}_i)_b)$. Its eigenvectors,
the principal axes, are the directions in which $\vb{L}$ is parallel to
$\boldsymbol{\omega}$.
\end{keyidea}

\notebook{05}{the inertia tensor of water and its principal axes}

\begin{selfcheck}
Two atoms of mass $m$ sit on the $z$ axis at $z = \pm d/2$. Write down
their inertia tensor about their midpoint, and its principal moments.
\end{selfcheck}
